% chap1.tex (Chapter 1 of the thesis)

\begin{document}

% Set page numbering to arabic the first time we commence a chapter.
% This is required to get the page numbering correct.
\pagenumbering{arabic}

% Note that the text in the [] brackets is the one that will
% appear in the table of contents, whilst the text in the {}
% brackets will appear in the main thesis.
\chapter[INTRODUCTION]{INTRODUCTION}

\section{Overview and introduction}

The primary aim of this thesis is to demonstrate a mechanism which can
be used to mediate changes in the mathematical or algorithmic
structure of a model of a complex system. A simple uptake-depuration
model of organisms which migrate through a contaminant field is
presented in Chapter \ref{cap2} to explore the fundamental principles
and to demonstrate the utility of such a mechanism. In this thesis, we
will use \emph{model} to refer to a (possibly trivial) collection of
of components which, together, simulate a system. The term
\emph{submodel} will generally refer to a largely discrete component
of a model, and it can be viewed as a model in its own right.  In some
sense we can view submodels as the abstract or \emph{in potentia}
representations, and \emph{agents} are specific, \emph{in esse}
instances of submodels.


Traditional models of complex systems incorporate alternative code
paths or expressions to deal with fundamentally different dynamics
which may arise during a simulaton. This is usually accomplished by
embedding specific tests which determine which code path or expression
is appropriate.  Agent based models are often (but not universally)
constructed with distinct kinds of agents having quite distinct code
bases {\citep{netlogo,swarm,repast,mason,nws,ningaloo}.


The models used as examples in this thesis
deal with such changes by changing the component used --- a number of
agents (instances of submodels)\footnote{}
associated with individual contact with a contaminant plume might be
substituted for a population agent which has no code path of its own
to deal with contaminant uptake or depuration.

Representing the state of the model, either as a whole or of its
constituent parts, is not simple: not only may the components of the
model vary through time, but their own dependences on other components
may change as the state of the model changes: a model which contains a
``whale spotting'' tourism venture may follow the activities of whales
at particular times of the year, but be utterly indifferent to them at
times when there is little likelihood of their presence in an
accessible location. Similarly, the association of entities
represented by a population-based model needs to be maintained if the
population disaggregates into super-individual-based or
individual-based agents.

These factors make a simple vector-based encoding mechanism for the
states of a model and its components awkward, thus we turn to a metric
space over a \rng\footnote{We use \rng\ to denote a ring without a
  multiplicative indentity element.} of trees with a finite number of
weighted, labelled nodes. This simplifies the comparison of possible
component mixes for a given global state, and makes available any
algorithms (such as for clustering) which depend only on ring
properties, excluding those that need a multiplicative unit.

Any attempt to assess the relative merit of a number of possible
component mixes for a model relative to its current state must have
some way of incorporating the necessary dependencies arising from the
effects of the current state on the , and be able to calculate (in
some sense) the incremental costs of change.


\section{Historical work}

Many individual-based models incorporate data which represents
environmental characteristics and influences the behaviour of of the
individuals simulated. \cite{Botkin72:2} and \cite{deangelis1978model}
are important early example which model individuals in in the context
of spatially explicit environmental gradients. \cite{Botkin72:2} In
both cases, the subjects of the simulations are represented as
individuals, and the environments are relatively static (though
interactions, such as shading in JABOWA, also played a role).  Even as
formulated, we can consider the environments in these models as static
models in their own right: from its early inception, the
individual-based modelling approach implicitly couples different
styles of model components.

The coastal marine ecosystem model in \cite{Gray06:1} used different
representations for the agents which represented the organisms
comprising the benthic habitat based on their life-stage; this nature
of this change was configured during the compilation of the model.
Juvenile biomasses were represented by polygonal clouds which could
change shape as they were advected. The an adult stage could be
represented by several different submodels, which either optimised for
speed or for spatial fidelity.  This model was in most other ways
similar to conventional agent-based modelling of the time.

%\citea{gross2002multimodeling} describes using this approach as a
%basic element underpinning their model, ATLSS, of the Everglades
%region in the U.S.A.

\cite{Bobashev2007hybrid} describes a model of epidemic simulation in
which the representation of populations or portions of populations are
decided based on the number of infected individuals relative to a
nominated trigger value. This model demonstrates that there is an
advantage to changeing representation in terms of computational
efficiency and the fidelity of the model. The model in Chapter
\ref{chap2} is similar: the rule governing the switching from one
representation to another occurs when some monitored quantity (the
number of infected individuals or the proximity to plume, for example)
crosses a nominated boundary.  The primary operational model
underpinning the work in \cite{Fulton2009crossingscales} takes a step
further, migrating whales change their representation according to
their putative geographic location, specifically whether they were in
the model domain or outside it. This model also had to consider
the linkages between agents: agents representing predators, wildlife
management and tourism operators had to be aware of the presence of
``in domain'' whales.


\section{Structure and function}










%% This new paragraph shows how to set \index{index items}index items
%% and \index{index items!subindex items} subindex items.

%% \subsection{New Subsection}
%% Here's a subsection with some simple maths $a^2+b^2=c^2$.

%% \subsubsection{subsubsection}
%% Here's a \index{subsubsection}subsubsection...oooooooohh....wow
%% wee!!!!!!

%% \newpage
%% Some more text to check indent and show how references work
%% \cite{Williamson:STS}.

%% Here's how we place a figure (Figure \ref{fig:utas}) on the page.
%% \begin{figure}[hbtp]
%% \begin{center}
%% \scalebox{.30}{\includegraphics{Utas_vert_BW}} \caption{
%% \label{fig:utas} The UTas logo}
%% \end{center}
%% \end{figure}

%% And finally, here's a table example (Table \ref{tab:taba}).
%% \begin{table}[hbtp]
%% \begin{center}
%% \begin{tabular}{|r|r|r|r|r|}
%% \hline
%% $n=$&2&3&4&5\\
%% \hline
%% $c$ (rad/day)&1.67&0.52&0.06&-0.17\\
%% \hline
%% period (days)&3.75&12.00&100.00&37.50\\
%% \hline
%% \end{tabular}
%% \end{center}
%% \caption{\label{tab:taba}A simple table}
%% \end{table}
